\scanpage{19}
% 1–5
\uline{Th. de préparation ``formel''} Soit $A$ un anneau linéairement topologisé, séparé et complet, $Z$ une indéterminée,
\[
F\deleted{(\unsure{Z})}=\sum_{i\geq 0}a_i Z^i\in A[[Z]],\quad N\in\mathbb N\ \text{tels que}
\]
(i)\qquad $a_i$ est top.$^{\mathrm t}$ nilpotent pour $0\leq i\leq N-1$

(ii)\qquad $a_N$ est inversible

% 6–10
Alors $F$ est un élément régulier de $A[[Z]]$ et $A[[Z]]/F$ a comme base $(Z^i)_{0\leq i\leq N-1}$, en d'autres termes tout élément $G$ de $A[[Z]]$ s'écrit de façon unique sous la forme
\[
\begin{aligned}
G&=QF+\sum_{0\leq i\leq N-1}b_iZ^i & Q&\in A[[Z]]\\
&&b_i&\in A
\end{aligned}
\]

% 11–15
\uline{Démonstration} Comme $A$ est séparé et complet pour sa top. linéaire, on est ramené aussitôt au cas où $A$ est discret, alors (i) signifie que les $a_i$ $(0\leq i\leq N-1)$ sont top. nilpotents. Soit $\mathfrak m$ l'idéal \inserted{(nilpotent)} qu'ils engendrent, $A_0=A/\mathfrak m$, et désignons par $\dot a_i,\dot F$ etc les réductions mod $\mathfrak m$.

% 16–22
\deleted{\unsure{Alors}} Donc on a \deleted{$\unsure{\dot a}$}
\[
\dot a_i=0\quad\text{pour}\quad 0\leq i\leq N-1
\]
i.e
\[
\begin{aligned}
\dot F&=\dot a_{N\deleted{\unsure{+1}}}Z^{N\deleted{\unsure{+1}}}+\ldots\\
&=\dot a_{N\deleted{+1}}Z^{N\deleted{+1}}(1+\ldots)
\end{aligned}
\]
donc, comme $\dot a_{N\deleted{\unsure{+1}}}$ est inversible par (ii), on voit que $\dot F$ est de la forme
\[
\dot F=\dot U Z^N,
\]
\deleted{\unsure{donc les|de sorte que les}} donc
\[
\unsure{\dot A_0|A_0}[[Z]]/\dot F\simeq A_0[[Z]]/Z^N\deleted{{}=A_0[Z]/Z^{\unsure{n|N}}}
\]
a comme base les $Z^i$ $(0\leq i\leq N-1)$

% 23–31
\deleted{i.e. (i) les $Z^i$ $(0\leq i\leq N-1)$ engendrent $A[[Z]]/FA[[Z]]$}

\deleted{En effet, ils engendrent mod $\mathfrak m$, on gagne donc par Nakayama.}

\deleted{\deleted{\unsure{Pour ce qui précède|Pour ce point il suffit de}}}

\deleted{(ii) $F$ est régulier. En effet, \deleted{\unsure{par récurrence sur l'ordre}} \inserted{le $\unsure{A\text{-module}|A_0\text{-module}}$ $A[[Z]]$ filtré $\mathfrak m$-adiquement} a un gr. associé isomorphe à \deleted{$\mathfrak m^n/\mathfrak m$} $\mathrm{Gr}^{*}(A)^{\mathbb N}$. \deleted{de nilpotence de $\mathfrak m$, SPS $\mathfrak m^2=0$. Soit $Q\in A[[Z]]/$ t.q. que $QF=0$, donc $Q\cdot$} et l'opération de multiplication \deleted{de $F$} par $F$ dessus, sur le gradué, correspond à : l'op. de multiplication par $\dot F$, qui est injective, OK.}
% archive sheet number: 18
